% General Measure Theory — Evans & Gariepy, Chapter 1 (pp. 1--59)
% Lecture notes style, Fall 2024, Mon/Wed schedule
% \input'd by main.tex


\section{August 26, 2024}\label{sec:lec1}

Pop quiz: what's a measure? If you said ``a function on sets,'' you're
not wrong but you're not really right either. The whole game in this
chapter is that you \emph{can't} define a measure on all subsets of $\R^n$
and keep the properties you want. So you start with something weaker ---
an outer measure --- and then cut down to the sets that behave.

Today we set up outer measures and Carath\'eodory's criterion.

\subsection{Outer measures}\label{subsec:outer-measures}

The idea is: outer measures are cheap. You can define them on
\emph{every} subset. The cost is you only get subadditivity, not
full additivity.

\begin{definition}[Outer measure]\label{def:outer-measure}
An \textbf{outer measure} on a set $X$ is a function
$\mu^* \colon \mathcal{P}(X) \to [0,\infty]$ satisfying
\begin{enumerate}[nosep]
  \item $\mu^*(\varnothing) = 0$,
  \item (Monotonicity) $A \subseteq B \implies \mu^*(A) \le \mu^*(B)$,
  \item (Countable subadditivity)
        $\mu^*\!\bigl(\bigcup_{k=1}^\infty A_k\bigr)
        \le \sum_{k=1}^\infty \mu^*(A_k)$.
\end{enumerate}
\end{definition}

That's it. No $\sigma$-algebra needed yet --- we're measuring everything, just
badly. The subadditivity is the price you pay for working with all of
$\mathcal{P}(X)$.

\subsection{Carath\'eodory's criterion}\label{subsec:caratheodory}

Here's where it gets interesting. We want to identify the ``well-behaved''
sets --- the ones where the outer measure actually splits additively.

\begin{definition}[$\mu^*$-measurability]\label{def:caratheodory}
A set $A \subseteq X$ is $\mu^*$-\textbf{measurable} if for every
$E \subseteq X$,
\[
  \mu^*(E) = \mu^*(E \cap A) + \mu^*(E \setminus A).
\]
\end{definition}

This definition confused me for a while. Why test against \emph{every} set $E$?
Brian explained it like this: you already get $\le$ from subadditivity for free.
So the real content is $\ge$. What you're asking is that $A$ acts as a
``clean cut'' --- when you slice any test set $E$ along $A$, no mass leaks out.
That's the whole point.

\begin{theorem}[Carath\'eodory]\label{thm:caratheodory}
Let $\mu^*$ be an outer measure on $X$. Then the collection
\[
  \mathcal{M} = \{ A \subseteq X : A \text{ is }
    \mu^*\text{-measurable} \}
\]
is a $\sigma$-algebra, and $\mu^*$ restricted to $\mathcal{M}$ is a
(complete) measure.
\end{theorem}

\begin{proof}[Proof sketch]
Complements are free --- the criterion is symmetric in $A$ and $A^c$.
For finite unions: take $A_1, A_2 \in \mathcal{M}$ and test against $E$.
Split first by $A_1$, then split $E \cap A_1^c$ by $A_2$. You get three
pieces that reassemble by subadditivity.

Now the hard part. Countable unions. Let $A = \bigcup_k A_k$ with the
$A_k$ disjoint (WLOG --- replace $A_k$ with
$A_k \setminus \bigcup_{j < k} A_j$, which stays in $\mathcal{M}$ by the
finite case). Set $B_n = \bigcup_{k=1}^n A_k$. Then for any $E$,
\[
  \mu^*(E \cap B_n) = \sum_{k=1}^n \mu^*(E \cap A_k)
\]
by induction. So
\[
  \mu^*(E) = \mu^*(E \cap B_n) + \mu^*(E \setminus B_n)
           \ge \sum_{k=1}^n \mu^*(E \cap A_k) + \mu^*(E \setminus A),
\]
since $A \supseteq B_n$. Let $n \to \infty$ and do the usual epsilon
management. Completeness is free: any $\mu^*$-null set passes the
criterion trivially. \qedhere
\end{proof}

\begin{remark}\label{rem:completeness-free}
That completeness comes for free is a big deal. If you build Lebesgue
measure via Carath\'eodory, it's automatically complete --- no need for
the separate completion step you'd do starting from Borel sets.
\end{remark}


\subsection{$\sigma$-algebras and Borel sets}\label{subsec:sigma-algebras}

\begin{definition}\label{def:sigma-algebra}
A $\sigma$-algebra on $X$ is a collection
$\mathcal{A} \subseteq \mathcal{P}(X)$ containing $\varnothing$, closed
under complements, and closed under countable unions.
\end{definition}

Nothing to see here. Closure under countable intersections follows from
De Morgan.

\begin{definition}[Borel $\sigma$-algebra]\label{def:borel}
For a topological space $X$, the \textbf{Borel $\sigma$-algebra}
$\mathcal{B}(X)$ is the $\sigma$-algebra generated by the open sets.
\end{definition}

In $\R^n$ you can equivalently generate $\mathcal{B}(\R^n)$ by open balls,
closed sets, compact sets, or half-spaces $\{x : x_i < a\}$. Evans doesn't
dwell on this. Every real analysis book covers it.


\subsection{Radon measures}\label{subsec:radon}

These regularity conditions are everywhere in later chapters, so I'm trying
to internalize the hierarchy now rather than looking it up every time.

\begin{definition}\label{def:borel-measure}
A \textbf{Borel measure} on $\R^n$ is a measure defined on
$\mathcal{B}(\R^n)$.
\end{definition}

\begin{definition}[Borel regular]\label{def:borel-regular}
A Borel measure $\mu$ is \textbf{Borel regular} if for every
$A \subseteq \R^n$, there exists a Borel set $B \supseteq A$ with
$\mu(B) = \mu^*(A)$.
\end{definition}

\begin{definition}[Radon measure]\label{def:radon}
A Borel regular measure $\mu$ on $\R^n$ is a \textbf{Radon measure} if
$\mu(K) < \infty$ for every compact $K \subseteq \R^n$.
\end{definition}

So the chain is: Radon = Borel regular + locally finite. Evans uses
``Radon'' as the default notion of ``good measure on $\R^n$.'' Lebesgue
measure $\mathcal{L}^n$ is the main example.

\begin{proposition}[Inner/outer regularity]\label{prop:radon-regularity}
If $\mu$ is a Radon measure on $\R^n$, then for every measurable set $A$:
\begin{enumerate}[nosep]
  \item (Outer regularity)
        $\mu(A) = \inf\{\mu(U) : U \supseteq A,\; U \text{ open}\}$.
  \item (Inner regularity)
        $\mu(A) = \sup\{\mu(K) : K \subseteq A,\; K \text{ compact}\}$.
\end{enumerate}
\end{proposition}

Proof: Evans pp.\ 5--7. Outer regularity uses Borel regularity. Inner
regularity uses $\sigma$-finiteness (exhaust by balls $B(0,k)$).

I spent two hours on the inner regularity proof before seeing the trick:
first prove it for finite-measure sets, then extend by writing
$A = \bigcup_k (A \cap B(0,k))$ and using continuity from below. This
took me a while.


\subsection{Measurable functions}\label{subsec:meas-fns}

\begin{definition}\label{def:measurable-fn}
Let $(X, \mathcal{A})$ be a measurable space. A function
$f \colon X \to [-\infty, \infty]$ is $\mathcal{A}$-\textbf{measurable}
if $\{f > t\} \in \mathcal{A}$ for all $t \in \R$.
\end{definition}

You could equivalently use $\{f \ge t\}$, $\{f < t\}$, or $\{f \le t\}$.
These are all the same --- you get from one to the other by complements and
countable intersections/unions. Evans uses the $\{f > t\}$ convention.

\begin{proposition}\label{prop:measurable-closure}
Measurable functions are closed under:
\begin{enumerate}[nosep]
  \item Pointwise limits: if $f_k \to f$ pointwise, each $f_k$ measurable,
        then $f$ is measurable.
  \item Sums, products, $\max$, $\min$, $|f|$.
  \item Countable $\sup$ and $\inf$: $\sup_k f_k$ and $\inf_k f_k$ are
        measurable.
\end{enumerate}
\end{proposition}

\begin{proof}[Proof of (iii)]
$\{\sup_k f_k > t\} = \bigcup_k \{f_k > t\}$. Countable union of
measurable sets. Done. The $\inf$ case is similar.
\qedhere
\end{proof}

\begin{remark}\label{rem:uncountable-sup}
Uncountable suprema can fail to be measurable. That's why pointwise
limits work (sequential = countable) but arbitrary sups don't.
Measurability is a countable-operations theory.
\end{remark}

TODO: Read Evans pp.\ 7--14 more carefully for the approximation by simple functions stuff.


% ===========================================================================
\section{August 28, 2024}\label{sec:lec2}

Question for the audience: if a sequence of measurable functions converges
pointwise, is the convergence ``almost uniform''? Turns out: yes, as long
as you're on a finite measure space. That's Egoroff.

Today we do Egoroff and Lusin, then start on integration.

\subsection{Egoroff's theorem}\label{subsec:egoroff}

This says convergence of measurable functions is ``almost uniform.'' The
proof is one of those things that looks like pure epsilon management but
the idea is actually pretty clean.

\begin{theorem}[Egoroff]\label{thm:egoroff}
Let $\mu$ be a finite measure on $X$, and let $f_k, f$ be measurable
functions with $f_k \to f$ a.e. Then for every $\varepsilon > 0$, there
exists a measurable set $A$ with $\mu(X \setminus A) < \varepsilon$ such
that $f_k \to f$ \textbf{uniformly} on $A$.
\end{theorem}

\begin{proof}
For each $m, n \in \N$, define
\[
  E_{m,n} = \bigcup_{k=n}^{\infty}
    \bigl\{ x : |f_k(x) - f(x)| \ge \tfrac{1}{m} \bigr\}.
\]
Then $E_{m,n} \downarrow$ as $n \to \infty$, and the a.e.\ convergence
gives $\mu\bigl(\bigcap_n E_{m,n}\bigr) = 0$. So by continuity from above
(here we use $\mu(X) < \infty$!), for each $m$ pick $n_m$ with
$\mu(E_{m,n_m}) < \varepsilon / 2^m$. Set
$A = X \setminus \bigcup_m E_{m,n_m}$. Then
$\mu(X \setminus A) < \varepsilon$, and on $A$ we have
$|f_k - f| < 1/m$ for all $k \ge n_m$. That's uniform convergence.
\qedhere
\end{proof}

The finiteness of $\mu$ is doing real work. Counterexample:
$f_k = \mathbf{1}_{[k, k+1]}$ on $\R$ with Lebesgue measure. Converges
to $0$ pointwise everywhere. But not uniformly on any set of finite
complement.

TODO: check if $\sigma$-finiteness is enough, or if we really need
$\mu(X) < \infty$. I think you need finiteness on the set where convergence
happens.


\subsection{Lusin's theorem}\label{subsec:lusin}

\begin{theorem}[Lusin]\label{thm:lusin}
Let $\mu$ be a Radon measure on $\R^n$ and $f \colon \R^n \to \R$
measurable. Then for every $\varepsilon > 0$ and every measurable set $A$
with $\mu(A) < \infty$, there exists a compact set $K \subseteq A$ with
$\mu(A \setminus K) < \varepsilon$ such that $f|_K$ is continuous.
\end{theorem}

Proof: Evans pp.\ 15--16. The idea is to use inner regularity to
approximate level sets $\{a_i \le f < a_{i+1}\}$ by compact sets. On the
union of those compact sets, $f$ is a staircase that can be refined to
continuous. Tietze extension shows up at one point.

Lusin is philosophically nice --- measurable functions are ``almost
continuous'' --- but I haven't needed it as much as Egoroff in practice.
It comes back for the Riesz representation theorem later.


\subsection{The Lebesgue integral}\label{subsec:integral}

Evans builds the integral the standard way.

\begin{definition}\label{def:integral}
For $f \colon X \to [0, \infty]$ measurable,
\[
  \int_X f \dd\mu = \sup\left\{ \int_X s \dd\mu :
    0 \le s \le f,\; s \text{ simple} \right\},
\]
where for a simple function $s = \sum_{i=1}^N a_i \mathbf{1}_{A_i}$
(with the $A_i$ measurable, disjoint),
$\int_X s \dd\mu = \sum_{i=1}^N a_i \,\mu(A_i)$.
\end{definition}

For general measurable $f$, write $f = f^+ - f^-$ and define
$\int f \dd\mu = \int f^+ \dd\mu - \int f^- \dd\mu$, provided at least one
side is finite. If both are finite, $f$ is \textbf{integrable}
($f \in L^1$). Nothing fancy.


\subsection{Monotone convergence theorem}\label{subsec:mct}

This is the foundation. DCT and Fatou both come from this.

\begin{theorem}[Monotone Convergence]\label{thm:mct}
Let $0 \le f_1 \le f_2 \le \cdots$ be measurable with
$f_k \nearrow f$ pointwise. Then
\[
  \int_X f \dd\mu = \lim_{k \to \infty} \int_X f_k \dd\mu.
\]
\end{theorem}

\begin{proof}
The $\le$ direction: $f_k \le f$ implies $\int f_k \le \int f$, so
$\lim \int f_k \le \int f$.

For $\ge$: pick any simple $s$ with $0 \le s \le f$ and fix $0 < t < 1$.
Define $A_k = \{x : f_k(x) \ge t \cdot s(x)\}$. Then $A_k \nearrow X$
(since $f_k \nearrow f \ge s > ts$ wherever $s > 0$). So
\[
  \int_X f_k \dd\mu \ge \int_{A_k} f_k \dd\mu
    \ge t \int_{A_k} s \dd\mu.
\]
Since $s$ is simple, $\int_{A_k} s \dd\mu \to \int_X s \dd\mu$ by
continuity from below. So $\lim \int f_k \ge t \int s$. Let
$t \nearrow 1$, then sup over simple $s \le f$. \qedhere
\end{proof}

I initially tried to use DCT to prove this. Circular. MCT has to come
first --- it's the foundation that everything else rests on.


\subsection{Fatou's lemma}\label{subsec:fatou}

\begin{lemma}[Fatou]\label{lem:fatou}
If $f_k \ge 0$ are measurable, then
\[
  \int_X \liminf_{k \to \infty} f_k \dd\mu
    \le \liminf_{k \to \infty} \int_X f_k \dd\mu.
\]
\end{lemma}

\begin{proof}
Set $g_k = \inf_{j \ge k} f_j$. Then $g_k \nearrow \liminf f_k$
and $g_k \le f_k$, so $\int g_k \le \int f_k$. Apply MCT to the left
side. \qedhere
\end{proof}

Short and clean. The inequality can be strict --- $f_k = k \cdot \mathbf{1}_{(0,1/k)}$
on $\R$: $\int f_k = 1$ for all $k$, but $\liminf f_k = 0$ a.e.


\subsection{Dominated convergence}\label{subsec:dct}

\begin{theorem}[Dominated Convergence]\label{thm:dct}
Let $f_k \to f$ a.e., with $|f_k| \le g$ a.e.\ for some integrable $g$.
Then $f$ is integrable and
\[
  \lim_{k \to \infty} \int_X f_k \dd\mu = \int_X f \dd\mu.
\]
\end{theorem}

\begin{proof}
The trick is to apply Fatou twice. $g + f_k \ge 0$ and $g - f_k \ge 0$, so:
\[
  \int (g + f) \le \liminf \int (g + f_k), \qquad
  \int (g - f) \le \liminf \int (g - f_k).
\]
Since $\int g < \infty$, cancel:
\[
  \int f \le \liminf \int f_k, \qquad
  -\int f \le -\limsup \int f_k.
\]
So $\limsup \int f_k \le \int f \le \liminf \int f_k$. We're done.
\qedhere
\end{proof}

The ``apply Fatou twice'' move is one of those things that's obvious after
you've seen it and baffling before. The domination hypothesis $|f_k| \le g$
is doing all the heavy lifting: it makes the Fatou inequalities go the right
way and keeps everything integrable.

\begin{remark}\label{rem:dct-counterex}
You can't drop the domination. Same counterexample:
$f_k = k \cdot \mathbf{1}_{(0,1/k)}$ converges to $0$ pointwise but
$\int f_k = 1$. No integrable dominator exists.
\end{remark}


% ===========================================================================
\section{September 4, 2024}\label{sec:lec3}

I accidentally tried to go to class on Monday. Labor Day. Nobody told
me.

Today: product measures, Fubini--Tonelli, and then the covering theorems
that set up the second half of the chapter.

\subsection{Product measures and Fubini's theorem}\label{subsec:product}

\begin{definition}[Product $\sigma$-algebra]\label{def:product-sigma}
If $(X, \mathcal{A})$ and $(Y, \mathcal{B})$ are measurable spaces, the
\textbf{product $\sigma$-algebra} $\mathcal{A} \otimes \mathcal{B}$ on
$X \times Y$ is generated by measurable rectangles $A \times B$
($A \in \mathcal{A}$, $B \in \mathcal{B}$).
\end{definition}

\begin{theorem}[Product measure]\label{thm:product-measure}
If $\mu$ and $\nu$ are $\sigma$-finite measures, there exists a unique
measure $\mu \otimes \nu$ on $\mathcal{A} \otimes \mathcal{B}$ with
$(\mu \otimes \nu)(A \times B) = \mu(A) \cdot \nu(B)$.
\end{theorem}

$\sigma$-finiteness is needed for uniqueness. Without it you can cook up
weird examples.

\begin{theorem}[Fubini--Tonelli]\label{thm:fubini}
Let $\mu, \nu$ be $\sigma$-finite measures.
\begin{enumerate}[nosep]
  \item \textbf{(Tonelli)} If $f \ge 0$ is
        $\mathcal{A} \otimes \mathcal{B}$-measurable, then
        \[
          \int_{X \times Y} f \dd(\mu \otimes \nu)
          = \int_X \!\left(\int_Y f(x,y) \dd\nu(y)\right) \dd\mu(x)
          = \int_Y \!\left(\int_X f(x,y) \dd\mu(x)\right) \dd\nu(y).
        \]
  \item \textbf{(Fubini)} If $f \in L^1(\mu \otimes \nu)$, the same
        holds, and the inner integrals are finite for a.e.\ value of the
        outer variable.
\end{enumerate}
\end{theorem}

The proof is a slog but the strategy is clean: verify for indicators of
measurable rectangles (trivial), extend to simple functions by linearity,
extend to nonneg measurable by MCT, extend to integrable by splitting
$f = f^+ - f^-$.

The point is that Tonelli is for $f \ge 0$ and doesn't need integrability,
while Fubini needs $f \in L^1$. Common trick: use Tonelli on $|f|$ to
check integrability, then use Fubini on $f$.


\subsubsection{Lebesgue measure}\label{subsubsec:lebesgue}

\begin{definition}\label{def:lebesgue-measure}
$n$-dimensional \textbf{Lebesgue measure} $\mathcal{L}^n$ is the product
$\underbrace{\mathcal{L}^1 \otimes \cdots \otimes \mathcal{L}^1}_{n}$,
where $\mathcal{L}^1$ is the completion of the Borel measure on $\R$
assigning to each interval $[a,b]$ its length $b-a$.
\end{definition}

Evans actually constructs $\mathcal{L}^n$ via the outer measure
\[
  \mathcal{L}^n_*(A)
  = \inf\!\left\{ \sum_{k=1}^\infty |Q_k| :
    A \subseteq \bigcup_{k=1}^\infty Q_k,\; Q_k \text{ cubes}\right\},
\]
where $|Q|$ is the volume of the cube. Carath\'eodory then gives the
$\sigma$-algebra and measure. The two constructions agree on Borel sets.

\begin{fact}\label{fact:lebesgue-props}
$\mathcal{L}^n$ is:
\begin{enumerate}[nosep]
  \item A Radon measure.
  \item Translation invariant:
        $\mathcal{L}^n(A + x) = \mathcal{L}^n(A)$.
  \item Scales correctly:
        $\mathcal{L}^n(rA) = r^n \mathcal{L}^n(A)$ for $r > 0$.
  \item Uniquely characterized (up to scalar) by translation invariance
        on Borel sets.
\end{enumerate}
\end{fact}


\subsection{Vitali's covering theorem}\label{subsec:vitali}

OK now we get into the technical stuff. Covering theorems aren't glamorous
but you literally can't differentiate measures without them. I found this
section painful on first reading.

\begin{theorem}[Vitali]\label{thm:vitali}
Let $\mathcal{F}$ be a collection of closed balls in $\R^n$ with
$\sup\{\diam(B) : B \in \mathcal{F}\} < \infty$. Then there exists a
countable subcollection of \textbf{disjoint} balls
$\{B_i\}_{i=1}^\infty \subseteq \mathcal{F}$ such that
\[
  \bigcup_{B \in \mathcal{F}} B
    \subseteq \bigcup_{i=1}^\infty \hat{B}_i,
\]
where $\hat{B}_i$ is the ball with the same center as $B_i$ but $5$
times the radius.
\end{theorem}

The ``$5r$-trick.'' The proof is a greedy algorithm: pick the biggest ball
(or close to it), throw out everything it intersects, repeat. Why $5$?
Because if two balls overlap and one has radius at most the other's, the
smaller one sits inside the $5\times$ enlargement. Nothing deeper than
that.

There's also a refined version for Radon measures:

\begin{theorem}[Vitali covering theorem for Radon measures]%
\label{thm:vitali-radon}
Let $\mu$ be a Radon measure, $A \subseteq \R^n$, and $\mathcal{F}$ a
collection of closed balls such that for each $x \in A$, there exist
balls in $\mathcal{F}$ centered at $x$ with arbitrarily small radius (a
``fine cover''). Then there exist disjoint balls
$\{B_i\} \subseteq \mathcal{F}$ with
\[
  \mu\!\left(A \setminus \bigcup_{i=1}^\infty B_i\right) = 0.
\]
\end{theorem}

So a fine Vitali cover lets you cover $\mu$-almost all of $A$ with
disjoint balls. Powerful.


\subsection{Besicovitch's covering theorem}\label{subsec:besicovitch}

\begin{theorem}[Besicovitch]\label{thm:besicovitch}
There exists a constant $N_n$ (depending only on dimension $n$) such
that: if $A \subseteq \R^n$ is bounded and $\mathcal{F}$ is a collection
of closed balls with a ball centered at each point of $A$, then there
exist at most $N_n$ countable subcollections
$\mathcal{G}_1, \ldots, \mathcal{G}_{N_n}$, each consisting of disjoint
balls, such that
\[
  A \subseteq \bigcup_{j=1}^{N_n}
    \bigcup_{B \in \mathcal{G}_j} B.
\]
\end{theorem}

Besicovitch is harder than Vitali. The proof is genuinely involved ---
Evans devotes pp.\ 30--35 to it. The advantage over Vitali: the balls
don't get enlarged. The trade-off: you get bounded overlap ($N_n$
families) instead of disjointness in one family. The constant $N_n$
depends only on dimension, and its exact value doesn't matter.

The geometric core of the proof is a counting argument: at any point,
the number of balls of comparable radius containing that point is bounded
by a dimensional constant. The proof is annoying but mechanical.

I won't reproduce it. I followed it once in lecture and I don't think I
could reproduce it cold.

TODO: redo the Besicovitch proof more carefully.


% ===========================================================================
\section{September 9, 2024}\label{sec:lec4}

Evans derives Radon--Nikodym as a \emph{consequence} of measure
differentiation, which is the reverse of the usual functional-analysis
approach. I think it's more geometric and more natural for what follows
(BV functions, sets of finite perimeter), but it makes this part of the
chapter harder. Worth it though.

\subsection{Differentiation of Radon measures}\label{subsec:differentiation}

This is where the covering theorems pay off. The goal: given two Radon
measures $\mu$ and $\nu$, differentiate $\nu$ with respect to $\mu$ by
looking at ratios on shrinking balls.

\begin{definition}\label{def:measure-deriv}
For Radon measures $\mu, \nu$ on $\R^n$, define
\[
  D_\mu \nu(x) = \lim_{r \to 0} \frac{\nu(B(x,r))}{\mu(B(x,r))},
\]
when the limit exists (and $\mu(B(x,r)) > 0$ for small $r$).
\end{definition}

Here's what makes this all work:

\begin{theorem}[Lebesgue--Besicovitch Differentiation]%
\label{thm:lebesgue-besicovitch}
Let $\mu$ be a Radon measure on $\R^n$ and $\nu$ a (signed) Radon
measure. Then:
\begin{enumerate}[nosep]
  \item $D_\mu \nu(x)$ exists and is finite $\mu$-a.e.
  \item The Lebesgue decomposition $\nu = \nu_{ac} + \nu_s$ (where
        $\nu_{ac} \ll \mu$ and $\nu_s \perp \mu$) satisfies
        $D_\mu \nu(x) = \frac{d\nu_{ac}}{d\mu}(x)$ \; $\mu$-a.e.
  \item $D_\mu \nu_s(x) = 0$ for $\mu$-a.e.\ $x$.
\end{enumerate}
\end{theorem}

Part (iii) is saying the singular part is ``invisible'' to the derivative.
I still find this proof non-obvious even after going through it twice.

\begin{proof}[Proof sketch]
Evans constructs the Lebesgue decomposition \emph{using} the derivative.
Define $D^+ = \limsup_{r \to 0}$ and $D_- = \liminf_{r \to 0}$ of the
ratio $\nu(B(x,r))/\mu(B(x,r))$. The main technical step: show that
\[
  \mu\{x : D^+ \nu(x) > t > s > D_- \nu(x)\} = 0
\]
for all $t > s$. This uses Besicovitch to build efficient covers at two
different scales. Once you know $D^+ = D_-$ a.e., the limit exists. The
identification with the Radon--Nikodym derivative is then a uniqueness
argument. See Evans pp.\ 39--43. \qedhere
\end{proof}

\begin{corollary}[Lebesgue decomposition]\label{cor:lebesgue-decomp}
Every signed Radon measure $\nu$ has a unique decomposition
$\nu = \nu_{ac} + \nu_s$ with $\nu_{ac} \ll \mu$ and $\nu_s \perp \mu$.
\end{corollary}

\begin{corollary}[Radon--Nikodym]\label{cor:radon-nikodym}
If $\nu \ll \mu$ (both Radon on $\R^n$), then there exists
$f \in L^1_{\mathrm{loc}}(\mu)$ with $\nu = f\mu$, i.e.,
$\nu(A) = \int_A f \dd\mu$ for all Borel $A$. And $f = D_\mu \nu$
a.e.
\end{corollary}


\subsection{Lebesgue points}\label{subsec:lebesgue-pts}

\begin{definition}[Lebesgue point]\label{def:lebesgue-point}
Let $f \in L^1_{\mathrm{loc}}(\R^n)$. A point $x$ is a
\textbf{Lebesgue point} of $f$ if
\[
  \lim_{r \to 0} \frac{1}{\mathcal{L}^n(B(x,r))}
    \int_{B(x,r)} |f(y) - f(x)| \dd y = 0.
\]
\end{definition}

This is stronger than saying the averages converge to $f(x)$. It's saying
the $L^1$ oscillation around $x$ vanishes. The distinction matters.

\begin{theorem}\label{thm:lebesgue-points}
If $f \in L^1_{\mathrm{loc}}(\R^n)$, then $\mathcal{L}^n$-a.e.\ $x$ is
a Lebesgue point of $f$.
\end{theorem}

\begin{proof}[Proof idea]
Apply Lebesgue--Besicovitch (Theorem~\ref{thm:lebesgue-besicovitch}) to
$\nu(A) = \int_A |f(y) - c| \dd y$ for each rational $c$. This gives
\[
  \lim_{r \to 0} \frac{1}{\mathcal{L}^n(B(x,r))} \int_{B(x,r)}
    |f(y) - c| \dd y = |f(x) - c| \quad \text{a.e.}
\]
for each rational $c$. Take a countable intersection (one for each
rational), then choose $c = c_k \to f(x)$ along rationals. You get the
full Lebesgue point condition. \qedhere
\end{proof}


\subsection{Approximate continuity}\label{subsec:approx-cont}

\begin{definition}\label{def:approx-cont}
A measurable function $f$ is \textbf{approximately continuous} at $x$ if
there exists a measurable set $A$ with density $1$ at $x$ (meaning
$\lim_{r \to 0}
\mathcal{L}^n(A \cap B(x,r)) / \mathcal{L}^n(B(x,r)) = 1$) such that
$f|_A$ is continuous at $x$.
\end{definition}

\begin{theorem}\label{thm:approx-cont}
If $f \in L^1_{\mathrm{loc}}(\R^n)$, then $f$ is approximately
continuous at $\mathcal{L}^n$-a.e.\ $x$.
\end{theorem}

This follows from the Lebesgue point theorem: if $x$ is a Lebesgue point,
take $A = \{y : |f(y) - f(x)| < \varepsilon\}$ and check it has density $1$.
Details on Evans p.\ 46.

These concepts --- Lebesgue points, approximate limits, approximate
continuity --- come back in a big way in Chapters~5 and~6 when Evans
discusses BV functions. For now I'm just filing them away.


% ===========================================================================
\section{September 11, 2024}\label{sec:lec5}

Last lecture on Chapter~1. We finish with the Riesz representation theorem
and weak$^*$ compactness for measures. This connects the measure theory to
functional analysis, and honestly this section feels the most ``structural''
compared to the covering-theorem grind from last week.

\subsection{Riesz representation theorem}\label{subsec:riesz}

This next result is why we built all the Radon measure machinery.

\begin{theorem}[Riesz Representation]\label{thm:riesz}
Let $L \colon C_c(\R^n) \to \R$ be a positive linear functional
(i.e., $L(f) \ge 0$ whenever $f \ge 0$). Then there exists a unique
Radon measure $\mu$ on $\R^n$ such that
\[
  L(f) = \int_{\R^n} f \dd\mu \quad
    \text{for all } f \in C_c(\R^n).
\]
\end{theorem}

Here $C_c(\R^n)$ is continuous functions with compact support.

\begin{proof}[Proof outline]
\textbf{Step 1.} Define $\mu$ on open sets by
\[
  \mu(U) = \sup\{ L(f) : f \in C_c(\R^n),\; 0 \le f \le 1,\;
    \spt(f) \subseteq U \}.
\]
\textbf{Step 2.} Extend to an outer measure via
$\mu^*(A) = \inf\{ \mu(U) : U \supseteq A,\; U \text{ open} \}$.

\textbf{Step 3.} Carath\'eodory gives a measure. Show it's Radon.

\textbf{Step 4.} Show $L(f) = \int f \dd\mu$ by approximating $f$ from
below with simple functions built from the sets $\{f > t\}$.

\textbf{Step 5.} Uniqueness: if two Radon measures agree on $C_c$, they
agree on open sets by inner regularity, hence on all Borel sets.
\end{proof}

I won't pretend I've internalized every detail. The construction in Step~1
is the key idea --- the linear functional ``tells you'' what the measure of
each open set should be, and you bootstrap from there. Actually, let me
back up. The reason this works is that open sets are determined from below
by continuous functions with compact support. That's the bridge.

TODO: work through the proof more carefully, maybe try it for $\R^1$
first.

\begin{remark}\label{rem:riesz-signed}
There's a signed version: every bounded linear functional on $C_0(\R^n)$
(continuous functions vanishing at infinity, with the sup norm) is
represented by a signed Radon measure. Evans doesn't state this but it's
standard (Rudin's \emph{Real and Complex Analysis}, Chapter~6).
\end{remark}


\subsection{Weak$^*$ convergence}\label{subsec:weak-star}

\begin{definition}\label{def:weak-star}
A sequence of Radon measures $\mu_k$ converges \textbf{weakly$^*$} to
$\mu$ (written $\mu_k \stackrel{*}{\rightharpoonup} \mu$) if
\[
  \int_{\R^n} f \dd\mu_k \to \int_{\R^n} f \dd\mu
  \quad \text{for all } f \in C_c(\R^n).
\]
\end{definition}

Evans calls this ``weak convergence'' but it's really weak$^*$ ---
measures live in $(C_c)^*$, not in a space that's itself a dual. This is
standard confusion in the literature and it used to trip me up.

\begin{theorem}[Weak$^*$ compactness]\label{thm:weak-compactness}
Let $\{\mu_k\}$ be Radon measures on $\R^n$ with
$\sup_k \mu_k(K) < \infty$ for every compact $K \subseteq \R^n$. Then
there exists a subsequence
$\mu_{k_j} \stackrel{*}{\rightharpoonup} \mu$ for some Radon measure
$\mu$.
\end{theorem}

\begin{proof}[Proof sketch]
$C_c(\R^n)$ is separable. Let $\{f_m\}$ be a countable dense subset. For
each $m$, the sequence $\{\int f_m \dd\mu_k\}$ is bounded by the uniform
bound on compact sets plus compact support of $f_m$. Diagonalize: extract
a subsequence so that $\int f_m \dd\mu_{k_j}$ converges for every $m$.

By density, $\int f \dd\mu_{k_j}$ converges for every
$f \in C_c$. This limit defines a positive linear functional, so by Riesz
(Theorem~\ref{thm:riesz}) it's integration against a Radon measure.
\qedhere
\end{proof}

So Riesz is doing real work here --- it converts the convergent functional
back into a measure.

\begin{proposition}[Lower semicontinuity]\label{prop:lsc-weak}
If $\mu_k \stackrel{*}{\rightharpoonup} \mu$ and
$U \subseteq \R^n$ is open, then
$\mu(U) \le \liminf_{k \to \infty} \mu_k(U)$.
\end{proposition}

\begin{proof}
Approximate $\mathbf{1}_U$ from below by $f \in C_c$ with $0 \le f \le 1$,
$\spt(f) \subseteq U$. Then
$\int f \dd\mu = \lim \int f \dd\mu_k \le \liminf \mu_k(U)$. Take sup
over such $f$. \qedhere
\end{proof}

This is one direction of the Portmanteau theorem. The other directions ---
$\limsup \mu_k(C) \le \mu(C)$ for closed $C$, convergence of $\mu_k(A)$ when
$\mu(\partial A) = 0$ --- follow by similar approximation.

\begin{example}\label{ex:weak-not-strong}
$\mu_k = \delta_{1/k}$ on $\R$. Then
$\mu_k \stackrel{*}{\rightharpoonup} \delta_0$ since $f(1/k) \to f(0)$
for all $f \in C_c(\R)$. But $\mu_k(\{0\}) = 0$ for all $k$ while
$\delta_0(\{0\}) = 1$. Weak convergence doesn't give convergence of
measures of specific sets unless the boundary has measure zero.
\end{example}

\begin{remark}\label{rem:prokhorov}
In probability, this compactness result is basically Prokhorov's theorem.
The condition $\sup_k \mu_k(K) < \infty$ is related to tightness. For
probability measures on $\R^n$, tightness and relative sequential
compactness in the weak topology are the same thing.
\end{remark}

\bigskip

That wraps up Chapter~1. A lot of machinery, but the payoff comes later
when Evans can freely differentiate measures, pass to weak limits, and
decompose into absolutely continuous and singular parts without reproof.
The covering theorems felt like the hardest part; the Riesz representation
and weak compactness feel the most useful going forward.

TODO: go back and do more examples for the differentiation section. I
understand Lebesgue--Besicovitch's statement but the proof still feels
fuzzy. Maybe try it for $\R^1$ where Besicovitch is basically just Vitali.
